\documentclass[11pt,a4paper]{article}
\newcommand{\gradelabel}{Grade 10}
\newcommand{\leconhdr}{Boucles et répétitions}
% Préambule commun aux plans de cours de robotique (Grades 7 à 12).
% Chaque plan de cours commence par \documentclass[11pt,a4paper]{article}, définit
% \gradelabel et \leconhdr (en-têtes de page), puis % Préambule commun aux plans de cours de robotique (Grades 7 à 12).
% Chaque plan de cours commence par \documentclass[11pt,a4paper]{article}, définit
% \gradelabel et \leconhdr (en-têtes de page), puis % Préambule commun aux plans de cours de robotique (Grades 7 à 12).
% Chaque plan de cours commence par \documentclass[11pt,a4paper]{article}, définit
% \gradelabel et \leconhdr (en-têtes de page), puis \input{../../preambule_robotique.tex}.
\usepackage[T1]{fontenc}
\usepackage[utf8]{inputenc}
\usepackage{lmodern}
\usepackage[french]{babel}
\usepackage[margin=2cm,headheight=14pt]{geometry}
\usepackage[most]{tcolorbox}
\usepackage{xcolor,tabularx,booktabs,array,enumitem,fancyhdr,amssymb,tikz}
\usetikzlibrary{arrows.meta,positioning}

\pagestyle{fancy}
\fancyhf{}
\lhead{\small Robotique -- \gradelabel{} -- Novembre 2026}
\rhead{\small \leconhdr}
\cfoot{\thepage}
\setlist{itemsep=1pt,topsep=2pt,leftmargin=*}
\renewcommand{\arraystretch}{1.35}
\newcolumntype{Y}{>{\raggedright\arraybackslash}X}
\newcommand{\cb}{$\square$\,}
\newcommand{\fl}{$\rightarrow$}
\newcommand{\lig}{\noindent\rule{\linewidth}{0.3pt}\\[5mm]}

\newtcolorbox{etape}[3]{enhanced,breakable,colback=#1!6,colframe=#1!70!black,
  coltitle=white,fonttitle=\bfseries,title={#2\hfill #3},boxrule=0.8pt,arc=2mm,
  before skip=8pt,after skip=8pt}
\newtcolorbox{astuce}[1][Différenciation et besoins particuliers]{breakable,
  colback=yellow!12,colframe=yellow!50!black,fonttitle=\bfseries\small,
  title={#1},boxrule=0.5pt,arc=1mm,before skip=6pt,after skip=6pt}
\newtcolorbox{code}[1][Code]{breakable,colback=black!4,colframe=black!60,
  fonttitle=\bfseries\small,title={#1},boxrule=0.5pt,arc=1mm,fontupper=\ttfamily\small,
  before skip=6pt,after skip=6pt}
% Blocs de code coloré (écrire le code sans accents : comme dans l'éditeur)
\newtcblisting{arduino}[1][Code Arduino]{listing only,breakable,colback=black!4,
  colframe=black!60,fonttitle=\bfseries\small,title={#1},boxrule=0.5pt,arc=1mm,
  before skip=6pt,after skip=6pt,
  listing options={language=C++,basicstyle=\ttfamily\small,columns=fullflexible,
  keepspaces=true,showstringspaces=false,commentstyle=\color{green!40!black}\itshape,
  keywordstyle=\color{blue!70!black}\bfseries}}
\newtcblisting{python}[1][Code Python]{listing only,breakable,colback=black!4,
  colframe=black!60,fonttitle=\bfseries\small,title={#1},boxrule=0.5pt,arc=1mm,
  before skip=6pt,after skip=6pt,
  listing options={language=Python,basicstyle=\ttfamily\small,columns=fullflexible,
  keepspaces=true,showstringspaces=false,commentstyle=\color{green!40!black}\itshape,
  keywordstyle=\color{blue!70!black}\bfseries}}
\newcommand{\entete}[2]{%
  \begin{tcolorbox}[colback=#1!10,colframe=#1!70!black,boxrule=0.8pt,arc=2mm]
  {\Large\bfseries #2}\hfill Nom : \rule{4cm}{0.3pt}\quad Date : \rule{2.2cm}{0.3pt}
  \end{tcolorbox}}

% Bandeau de titre : #1 titre, #2 semaine
\newcommand{\titreblock}[2]{%
\begin{tcolorbox}[colback=teal!10,colframe=teal!70!black,arc=2mm,boxrule=1pt]
\centering
{\LARGE\bfseries #1}\\[2mm]
{\large Robotique -- \gradelabel{} -- Plan de cours de 2 heures}\\[1mm]
Enseignant : Francis Ki \quad|\quad #2
\end{tcolorbox}}

% Tableau du minutage : 6 contenus (accroche, mini-leçon, guidée, autonome, mise en commun, exit ticket)
\newcommand{\minutage}[6]{%
\begin{tabularx}{\linewidth}{>{\bfseries}p{2.2cm}p{2.6cm}Y p{3cm}}
\toprule
Minutes & Étape & Contenu & Organisation\\
\midrule
0 -- 10 & Accroche & #1 & Classe entière\\
10 -- 30 & Mini-leçon & #2 & Classe entière\\
30 -- 60 & Pratique guidée & #3 & Collectif puis binômes\\
60 -- 95 & Pratique autonome & #4 & Individuel sur poste\\
95 -- 110 & Mise en commun & #5 & Classe entière\\
110 -- 120 & Exit ticket & #6 & Individuel sur poste\\
\bottomrule
\end{tabularx}}

% Étayage progressif : #1 aide 1, #2 aide 2, #3 aide 3
\newcommand{\aides}[3]{%
\textbf{Étayage progressif de l'aide.}
\begin{center}
\begin{tabularx}{\linewidth}{>{\bfseries}p{2.4cm}Y p{4.4cm}}
\toprule
Niveau d'aide & Contenu & Qui intervient\\
\midrule
Aide 1 & #1 & Le binôme seul\\
Aide 2 & #2 & Le binôme avec la carte d'aide\\
Aide 3 & #3 & L'enseignant\\
\bottomrule
\end{tabularx}
\end{center}}

% Les trois groupes de la pratique autonome : #1 tâche A, #2 tâche B, #3 tâche C
\newcommand{\groupes}[3]{%
\textbf{Consigne commune\,:} chaque élève ouvre le fichier de son groupe, le renomme avec son prénom et travaille seul sur son poste. L'enseignant passe d'abord auprès du groupe A pendant 10 min.

\medskip
\begin{tabularx}{\linewidth}{>{\bfseries}p{3cm}Y p{2.2cm}}
\toprule
Groupe & Tâche sur poste & Annexe\\
\midrule
A : en dessous du niveau & #1 & Annexe C\\
B : au niveau & #2 & Annexe D\\
C : au-dessus du niveau & #3 & Annexe E\\
\bottomrule
\end{tabularx}}

% Mise en commun : #1 groupe A, #2 groupe B, #3 groupe C, #4 conclusion
\newcommand{\miseencommun}[4]{%
\begin{etape}{teal}{5. Mise en commun}{95 -- 110 min}
\begin{enumerate}
 \item (2 min) Retour au calme, écrans mis en veille.
 \item (9 min) Trois passages courts de 3 min, avec valorisation :
 \begin{itemize}
  \item \textbf{Groupe A} : #1
  \item \textbf{Groupe B} : #2
  \item \textbf{Groupe C} : #3
 \end{itemize}
 \item (4 min) Conclusion : #4
\end{enumerate}
\end{etape}}

% Clôture : #1 question 1, #2 question 2
\newcommand{\cloture}[2]{%
\begin{etape}{gray}{6. Clôture : exit ticket}{110 -- 120 min}
Formulaire en ligne ou papier (annexe G) :
\begin{enumerate}
 \item #1
 \item #2
 \item Auto-évaluation : \emph{Je maîtrise / Je progresse / J'ai besoin d'aide.}
\end{enumerate}
Lire les réponses avant la séance suivante pour ajuster les groupes de besoin.
\end{etape}}

% Grille d'observation : #1 #2 #3 intitulés de colonnes
\newcommand{\grille}[3]{%
\subsection*{Grille d'observation (enseignant)}
\begin{tabularx}{\linewidth}{Y p{2.4cm} p{2.4cm} p{2.4cm} Y}
\toprule
Élève & #1 & #2 & #3 & Groupe suivant\\
\midrule
 & & & & \\ & & & & \\ & & & & \\ & & & & \\ & & & & \\
\bottomrule
\end{tabularx}}

% Signature de l'exit ticket
\newcommand{\autoeval}{Je me sens : \cb Je maîtrise \quad \cb Je progresse \quad \cb J'ai besoin d'aide}
.
\usepackage[T1]{fontenc}
\usepackage[utf8]{inputenc}
\usepackage{lmodern}
\usepackage[french]{babel}
\usepackage[margin=2cm,headheight=14pt]{geometry}
\usepackage[most]{tcolorbox}
\usepackage{xcolor,tabularx,booktabs,array,enumitem,fancyhdr,amssymb,tikz}
\usetikzlibrary{arrows.meta,positioning}

\pagestyle{fancy}
\fancyhf{}
\lhead{\small Robotique -- \gradelabel{} -- Novembre 2026}
\rhead{\small \leconhdr}
\cfoot{\thepage}
\setlist{itemsep=1pt,topsep=2pt,leftmargin=*}
\renewcommand{\arraystretch}{1.35}
\newcolumntype{Y}{>{\raggedright\arraybackslash}X}
\newcommand{\cb}{$\square$\,}
\newcommand{\fl}{$\rightarrow$}
\newcommand{\lig}{\noindent\rule{\linewidth}{0.3pt}\\[5mm]}

\newtcolorbox{etape}[3]{enhanced,breakable,colback=#1!6,colframe=#1!70!black,
  coltitle=white,fonttitle=\bfseries,title={#2\hfill #3},boxrule=0.8pt,arc=2mm,
  before skip=8pt,after skip=8pt}
\newtcolorbox{astuce}[1][Différenciation et besoins particuliers]{breakable,
  colback=yellow!12,colframe=yellow!50!black,fonttitle=\bfseries\small,
  title={#1},boxrule=0.5pt,arc=1mm,before skip=6pt,after skip=6pt}
\newtcolorbox{code}[1][Code]{breakable,colback=black!4,colframe=black!60,
  fonttitle=\bfseries\small,title={#1},boxrule=0.5pt,arc=1mm,fontupper=\ttfamily\small,
  before skip=6pt,after skip=6pt}
% Blocs de code coloré (écrire le code sans accents : comme dans l'éditeur)
\newtcblisting{arduino}[1][Code Arduino]{listing only,breakable,colback=black!4,
  colframe=black!60,fonttitle=\bfseries\small,title={#1},boxrule=0.5pt,arc=1mm,
  before skip=6pt,after skip=6pt,
  listing options={language=C++,basicstyle=\ttfamily\small,columns=fullflexible,
  keepspaces=true,showstringspaces=false,commentstyle=\color{green!40!black}\itshape,
  keywordstyle=\color{blue!70!black}\bfseries}}
\newtcblisting{python}[1][Code Python]{listing only,breakable,colback=black!4,
  colframe=black!60,fonttitle=\bfseries\small,title={#1},boxrule=0.5pt,arc=1mm,
  before skip=6pt,after skip=6pt,
  listing options={language=Python,basicstyle=\ttfamily\small,columns=fullflexible,
  keepspaces=true,showstringspaces=false,commentstyle=\color{green!40!black}\itshape,
  keywordstyle=\color{blue!70!black}\bfseries}}
\newcommand{\entete}[2]{%
  \begin{tcolorbox}[colback=#1!10,colframe=#1!70!black,boxrule=0.8pt,arc=2mm]
  {\Large\bfseries #2}\hfill Nom : \rule{4cm}{0.3pt}\quad Date : \rule{2.2cm}{0.3pt}
  \end{tcolorbox}}

% Bandeau de titre : #1 titre, #2 semaine
\newcommand{\titreblock}[2]{%
\begin{tcolorbox}[colback=teal!10,colframe=teal!70!black,arc=2mm,boxrule=1pt]
\centering
{\LARGE\bfseries #1}\\[2mm]
{\large Robotique -- \gradelabel{} -- Plan de cours de 2 heures}\\[1mm]
Enseignant : Francis Ki \quad|\quad #2
\end{tcolorbox}}

% Tableau du minutage : 6 contenus (accroche, mini-leçon, guidée, autonome, mise en commun, exit ticket)
\newcommand{\minutage}[6]{%
\begin{tabularx}{\linewidth}{>{\bfseries}p{2.2cm}p{2.6cm}Y p{3cm}}
\toprule
Minutes & Étape & Contenu & Organisation\\
\midrule
0 -- 10 & Accroche & #1 & Classe entière\\
10 -- 30 & Mini-leçon & #2 & Classe entière\\
30 -- 60 & Pratique guidée & #3 & Collectif puis binômes\\
60 -- 95 & Pratique autonome & #4 & Individuel sur poste\\
95 -- 110 & Mise en commun & #5 & Classe entière\\
110 -- 120 & Exit ticket & #6 & Individuel sur poste\\
\bottomrule
\end{tabularx}}

% Étayage progressif : #1 aide 1, #2 aide 2, #3 aide 3
\newcommand{\aides}[3]{%
\textbf{Étayage progressif de l'aide.}
\begin{center}
\begin{tabularx}{\linewidth}{>{\bfseries}p{2.4cm}Y p{4.4cm}}
\toprule
Niveau d'aide & Contenu & Qui intervient\\
\midrule
Aide 1 & #1 & Le binôme seul\\
Aide 2 & #2 & Le binôme avec la carte d'aide\\
Aide 3 & #3 & L'enseignant\\
\bottomrule
\end{tabularx}
\end{center}}

% Les trois groupes de la pratique autonome : #1 tâche A, #2 tâche B, #3 tâche C
\newcommand{\groupes}[3]{%
\textbf{Consigne commune\,:} chaque élève ouvre le fichier de son groupe, le renomme avec son prénom et travaille seul sur son poste. L'enseignant passe d'abord auprès du groupe A pendant 10 min.

\medskip
\begin{tabularx}{\linewidth}{>{\bfseries}p{3cm}Y p{2.2cm}}
\toprule
Groupe & Tâche sur poste & Annexe\\
\midrule
A : en dessous du niveau & #1 & Annexe C\\
B : au niveau & #2 & Annexe D\\
C : au-dessus du niveau & #3 & Annexe E\\
\bottomrule
\end{tabularx}}

% Mise en commun : #1 groupe A, #2 groupe B, #3 groupe C, #4 conclusion
\newcommand{\miseencommun}[4]{%
\begin{etape}{teal}{5. Mise en commun}{95 -- 110 min}
\begin{enumerate}
 \item (2 min) Retour au calme, écrans mis en veille.
 \item (9 min) Trois passages courts de 3 min, avec valorisation :
 \begin{itemize}
  \item \textbf{Groupe A} : #1
  \item \textbf{Groupe B} : #2
  \item \textbf{Groupe C} : #3
 \end{itemize}
 \item (4 min) Conclusion : #4
\end{enumerate}
\end{etape}}

% Clôture : #1 question 1, #2 question 2
\newcommand{\cloture}[2]{%
\begin{etape}{gray}{6. Clôture : exit ticket}{110 -- 120 min}
Formulaire en ligne ou papier (annexe G) :
\begin{enumerate}
 \item #1
 \item #2
 \item Auto-évaluation : \emph{Je maîtrise / Je progresse / J'ai besoin d'aide.}
\end{enumerate}
Lire les réponses avant la séance suivante pour ajuster les groupes de besoin.
\end{etape}}

% Grille d'observation : #1 #2 #3 intitulés de colonnes
\newcommand{\grille}[3]{%
\subsection*{Grille d'observation (enseignant)}
\begin{tabularx}{\linewidth}{Y p{2.4cm} p{2.4cm} p{2.4cm} Y}
\toprule
Élève & #1 & #2 & #3 & Groupe suivant\\
\midrule
 & & & & \\ & & & & \\ & & & & \\ & & & & \\ & & & & \\
\bottomrule
\end{tabularx}}

% Signature de l'exit ticket
\newcommand{\autoeval}{Je me sens : \cb Je maîtrise \quad \cb Je progresse \quad \cb J'ai besoin d'aide}
.
\usepackage[T1]{fontenc}
\usepackage[utf8]{inputenc}
\usepackage{lmodern}
\usepackage[french]{babel}
\usepackage[margin=2cm,headheight=14pt]{geometry}
\usepackage[most]{tcolorbox}
\usepackage{xcolor,tabularx,booktabs,array,enumitem,fancyhdr,amssymb,tikz}
\usetikzlibrary{arrows.meta,positioning}

\pagestyle{fancy}
\fancyhf{}
\lhead{\small Robotique -- \gradelabel{} -- Novembre 2026}
\rhead{\small \leconhdr}
\cfoot{\thepage}
\setlist{itemsep=1pt,topsep=2pt,leftmargin=*}
\renewcommand{\arraystretch}{1.35}
\newcolumntype{Y}{>{\raggedright\arraybackslash}X}
\newcommand{\cb}{$\square$\,}
\newcommand{\fl}{$\rightarrow$}
\newcommand{\lig}{\noindent\rule{\linewidth}{0.3pt}\\[5mm]}

\newtcolorbox{etape}[3]{enhanced,breakable,colback=#1!6,colframe=#1!70!black,
  coltitle=white,fonttitle=\bfseries,title={#2\hfill #3},boxrule=0.8pt,arc=2mm,
  before skip=8pt,after skip=8pt}
\newtcolorbox{astuce}[1][Différenciation et besoins particuliers]{breakable,
  colback=yellow!12,colframe=yellow!50!black,fonttitle=\bfseries\small,
  title={#1},boxrule=0.5pt,arc=1mm,before skip=6pt,after skip=6pt}
\newtcolorbox{code}[1][Code]{breakable,colback=black!4,colframe=black!60,
  fonttitle=\bfseries\small,title={#1},boxrule=0.5pt,arc=1mm,fontupper=\ttfamily\small,
  before skip=6pt,after skip=6pt}
% Blocs de code coloré (écrire le code sans accents : comme dans l'éditeur)
\newtcblisting{arduino}[1][Code Arduino]{listing only,breakable,colback=black!4,
  colframe=black!60,fonttitle=\bfseries\small,title={#1},boxrule=0.5pt,arc=1mm,
  before skip=6pt,after skip=6pt,
  listing options={language=C++,basicstyle=\ttfamily\small,columns=fullflexible,
  keepspaces=true,showstringspaces=false,commentstyle=\color{green!40!black}\itshape,
  keywordstyle=\color{blue!70!black}\bfseries}}
\newtcblisting{python}[1][Code Python]{listing only,breakable,colback=black!4,
  colframe=black!60,fonttitle=\bfseries\small,title={#1},boxrule=0.5pt,arc=1mm,
  before skip=6pt,after skip=6pt,
  listing options={language=Python,basicstyle=\ttfamily\small,columns=fullflexible,
  keepspaces=true,showstringspaces=false,commentstyle=\color{green!40!black}\itshape,
  keywordstyle=\color{blue!70!black}\bfseries}}
\newcommand{\entete}[2]{%
  \begin{tcolorbox}[colback=#1!10,colframe=#1!70!black,boxrule=0.8pt,arc=2mm]
  {\Large\bfseries #2}\hfill Nom : \rule{4cm}{0.3pt}\quad Date : \rule{2.2cm}{0.3pt}
  \end{tcolorbox}}

% Bandeau de titre : #1 titre, #2 semaine
\newcommand{\titreblock}[2]{%
\begin{tcolorbox}[colback=teal!10,colframe=teal!70!black,arc=2mm,boxrule=1pt]
\centering
{\LARGE\bfseries #1}\\[2mm]
{\large Robotique -- \gradelabel{} -- Plan de cours de 2 heures}\\[1mm]
Enseignant : Francis Ki \quad|\quad #2
\end{tcolorbox}}

% Tableau du minutage : 6 contenus (accroche, mini-leçon, guidée, autonome, mise en commun, exit ticket)
\newcommand{\minutage}[6]{%
\begin{tabularx}{\linewidth}{>{\bfseries}p{2.2cm}p{2.6cm}Y p{3cm}}
\toprule
Minutes & Étape & Contenu & Organisation\\
\midrule
0 -- 10 & Accroche & #1 & Classe entière\\
10 -- 30 & Mini-leçon & #2 & Classe entière\\
30 -- 60 & Pratique guidée & #3 & Collectif puis binômes\\
60 -- 95 & Pratique autonome & #4 & Individuel sur poste\\
95 -- 110 & Mise en commun & #5 & Classe entière\\
110 -- 120 & Exit ticket & #6 & Individuel sur poste\\
\bottomrule
\end{tabularx}}

% Étayage progressif : #1 aide 1, #2 aide 2, #3 aide 3
\newcommand{\aides}[3]{%
\textbf{Étayage progressif de l'aide.}
\begin{center}
\begin{tabularx}{\linewidth}{>{\bfseries}p{2.4cm}Y p{4.4cm}}
\toprule
Niveau d'aide & Contenu & Qui intervient\\
\midrule
Aide 1 & #1 & Le binôme seul\\
Aide 2 & #2 & Le binôme avec la carte d'aide\\
Aide 3 & #3 & L'enseignant\\
\bottomrule
\end{tabularx}
\end{center}}

% Les trois groupes de la pratique autonome : #1 tâche A, #2 tâche B, #3 tâche C
\newcommand{\groupes}[3]{%
\textbf{Consigne commune\,:} chaque élève ouvre le fichier de son groupe, le renomme avec son prénom et travaille seul sur son poste. L'enseignant passe d'abord auprès du groupe A pendant 10 min.

\medskip
\begin{tabularx}{\linewidth}{>{\bfseries}p{3cm}Y p{2.2cm}}
\toprule
Groupe & Tâche sur poste & Annexe\\
\midrule
A : en dessous du niveau & #1 & Annexe C\\
B : au niveau & #2 & Annexe D\\
C : au-dessus du niveau & #3 & Annexe E\\
\bottomrule
\end{tabularx}}

% Mise en commun : #1 groupe A, #2 groupe B, #3 groupe C, #4 conclusion
\newcommand{\miseencommun}[4]{%
\begin{etape}{teal}{5. Mise en commun}{95 -- 110 min}
\begin{enumerate}
 \item (2 min) Retour au calme, écrans mis en veille.
 \item (9 min) Trois passages courts de 3 min, avec valorisation :
 \begin{itemize}
  \item \textbf{Groupe A} : #1
  \item \textbf{Groupe B} : #2
  \item \textbf{Groupe C} : #3
 \end{itemize}
 \item (4 min) Conclusion : #4
\end{enumerate}
\end{etape}}

% Clôture : #1 question 1, #2 question 2
\newcommand{\cloture}[2]{%
\begin{etape}{gray}{6. Clôture : exit ticket}{110 -- 120 min}
Formulaire en ligne ou papier (annexe G) :
\begin{enumerate}
 \item #1
 \item #2
 \item Auto-évaluation : \emph{Je maîtrise / Je progresse / J'ai besoin d'aide.}
\end{enumerate}
Lire les réponses avant la séance suivante pour ajuster les groupes de besoin.
\end{etape}}

% Grille d'observation : #1 #2 #3 intitulés de colonnes
\newcommand{\grille}[3]{%
\subsection*{Grille d'observation (enseignant)}
\begin{tabularx}{\linewidth}{Y p{2.4cm} p{2.4cm} p{2.4cm} Y}
\toprule
Élève & #1 & #2 & #3 & Groupe suivant\\
\midrule
 & & & & \\ & & & & \\ & & & & \\ & & & & \\ & & & & \\
\bottomrule
\end{tabularx}}

% Signature de l'exit ticket
\newcommand{\autoeval}{Je me sens : \cb Je maîtrise \quad \cb Je progresse \quad \cb J'ai besoin d'aide}


\begin{document}
\titreblock{Boucles et répétitions (Python)}{Semaine du 16 au 20 novembre 2026}

%==================================================================
\section{Informations générales}
\begin{tabularx}{\linewidth}{>{\bfseries}p{3.6cm}Y}
\toprule
Titre & Boucles \texttt{for} et \texttt{while}\\
Classe / effectif & Grade 10, 11 à 20 élèves, niveaux hétérogènes, élèves à besoins particuliers\\
Durée & 2 heures (120 min)\\
Lieu & Salle informatique, un poste par élève\\
Matériel et logiciels &
\begin{itemize}
 \item Un poste par élève avec l'éditeur \textbf{Mu} en mode \emph{Python 3} (bouton Arrêter pour interrompre une boucle infinie, Débogueur pour suivre l'exécution).
 \item Vidéoprojecteur ou tableau interactif avec Mu projeté.
 \item Un dossier partagé contenant les fichiers de départ \texttt{.py} des trois niveaux.
 \item Inspiration : unité «~Round and round~» du cours TeachComputing (binômes, codage en direct, exemples commentés, idées fausses courantes).
\end{itemize}\\
Pré-requis &
Séances précédentes : variables, types, saisie, conditions \texttt{if / elif / else} et indentation.\\
\bottomrule
\end{tabularx}

\begin{astuce}[Préparation avant le cours (environ 40 min)]
\begin{enumerate}
 \item Préparer les fichiers de départ : \texttt{niveauA.py} (boucle for à compléter), \texttt{niveauB.py} (vide avec consignes), \texttt{niveauC.py} (début du jeu «~devine le nombre~»).
 \item Tester le bouton Arrêter de Mu avec une boucle infinie volontaire.
 \item Créer le formulaire de sortie (annexe G) ou imprimer l'annexe.
\end{enumerate}
\end{astuce}

\begin{astuce}[Prévenir la boucle infinie]
Montrer dès la mini-leçon ce qui se passe avec \texttt{while True} sans arrêt, et comment arrêter l'exécution dans Mu. Cela évite la panique en classe et sert de leçon sur l'erreur de logique.
\end{astuce}

%==================================================================
\section{Objectifs d'apprentissage}
\begin{tcolorbox}[colback=gray!8,colframe=gray!60!black,arc=1mm,boxrule=0.5pt]
\textbf{Objectifs du programme}
\begin{itemize}
 \item Utiliser les boucles \texttt{for} et \texttt{while}.
 \item Automatiser des traitements répétitifs.
\end{itemize}
\textbf{À la fin de la séance, l'élève est capable\,:}
\begin{itemize}
 \item d'écrire une boucle \texttt{for} avec \texttt{range} et de prévoir les valeurs parcourues\,;
 \item d'écrire une boucle \texttt{while} avec une condition d'arrêt et un compteur\,;
 \item de choisir entre \texttt{for} (nombre de tours connu) et \texttt{while} (jusqu'à une condition).
\end{itemize}
\textbf{Compétences numériques\,:} suivre l'exécution avec une table de trace, utiliser le débogueur, interrompre un programme.
\end{tcolorbox}

%==================================================================
\section{Déroulé minuté (120 min)}
\minutage{«~Je ne bavarderai plus~» cent fois}{\texttt{for}, \texttt{range}, \texttt{while}}{Une table de multiplication en collectif, puis un mot de passe en binômes}{Trois niveaux de besoin}{Valoriser chaque groupe}{Trois questions sur poste}

%------------------------------------------------------------------
\begin{etape}{orange}{1. Accroche : cent lignes de punition}{0 -- 10 min}
\begin{enumerate}
 \item (5 min) \emph{Un enseignant demande d'écrire cent fois «~Je ne bavarderai plus~». Comment un ordinateur le ferait-il sans que vous écriviez cent lignes ?} Les élèves proposent des idées.
 \item (5 min) Écrire deux versions au tableau : cent \texttt{print} recopiés, ou une ligne qui répète. Question-problème : \textbf{«~Comment dire à l'ordinateur de répéter une action ?~»}
\end{enumerate}
\end{etape}

%------------------------------------------------------------------
\begin{etape}{blue}{2. Mini-leçon : for et while}{10 -- 30 min}
\textbf{Support\,:} annexe A. Codage en direct. Quatre temps de 5 min.

\textbf{Temps 1 : for et range.}
\begin{python}[Boucle for]
for i in range(5):
    print("Tour numero", i)
\end{python}
\texttt{range(5)} donne 0, 1, 2, 3, 4 (cinq valeurs, on commence à 0). \texttt{range(1, 6)} donne 1 à 5. \texttt{range(0, 10, 2)} donne 0, 2, 4, 6, 8.

\textbf{Temps 2 : while.}
\begin{python}[Boucle while]
n = 3
while n > 0:
    print(n)
    n = n - 1
print("Decollage !")
\end{python}
La boucle se répète \emph{tant que} la condition est vraie. Il faut modifier la variable pour que la condition devienne fausse.

\textbf{Temps 3 : la table de trace.}
\begin{center}
\begin{tabularx}{0.8\linewidth}{>{\bfseries}c c c}
\toprule
Tour & Valeur de \texttt{n} avant & Affichage\\
\midrule
1 & 3 & 3\\
2 & 2 & 2\\
3 & 1 & 1\\
\bottomrule
\end{tabularx}
\end{center}

\textbf{Temps 4 : choisir.} \texttt{for} quand on connaît le nombre de tours, \texttt{while} quand on répète jusqu'à ce qu'une condition change (mot de passe, saisie valide, jeu).

\begin{astuce}
Faire suivre la table de trace avec un doigt ou un stylo. Utiliser le débogueur de Mu pour voir la variable changer ligne après ligne. Rappeler le piège classique : \texttt{range(5)} ne donne pas 5.
\end{astuce}
\end{etape}

%------------------------------------------------------------------
\begin{etape}{green}{3. Pratique guidée}{30 -- 60 min}
\textbf{Fiche\,:} annexe B. Programmation en binômes (pilote et navigateur, rôles échangés à mi-temps).

\textbf{3a. Table de multiplication en collectif (12 min).} Écrire en codage en direct un programme qui demande un nombre puis affiche sa table de 1 à 10 avec une boucle \texttt{for}.

\textbf{3b. Mot de passe en binômes (18 min).} Écrire un programme qui redemande un mot de passe tant que la réponse est fausse (\texttt{while}), puis affiche «~Acces autorise~».

\aides{Le programme du 3a et l'annexe A avec la table de trace}{Questions guides : \emph{Combien de fois la boucle se répète-t-elle ? Quelle variable change à chaque tour ? Quand la condition devient-elle fausse ?}}{L'enseignant écrit la ligne \texttt{while} avec le binôme, qui continue}

\begin{astuce}
Écrire d'abord le programme en français («~tant que le mot de passe est faux, redemander~»), puis le traduire en Python. Si la boucle ne s'arrête pas, cliquer sur Arrêter et chercher la ligne qui devrait modifier la variable.
\end{astuce}
\end{etape}

%------------------------------------------------------------------
\begin{etape}{violet}{4. Pratique autonome différenciée}{60 -- 95 min}
\groupes{Fichier \texttt{niveauA.py} à compléter : boucle \texttt{for} qui affiche les nombres de 1 à 10, puis table de trace à remplir, avec une banque de mots (\texttt{for}, \texttt{in}, \texttt{range}, \texttt{print}).}{Écrire trois programmes : table de multiplication choisie par l'utilisateur, somme des nombres de 1 à $n$ avec \texttt{for}, mot de passe à trois essais avec \texttt{while}. Tester chacun avec deux jeux d'essais.}{Réaliser le jeu «~devine le nombre~» : l'ordinateur choisit un nombre entre 1 et 100 (\texttt{random.randint}), le joueur propose, le programme indique «~plus grand~» ou «~plus petit~», compte les essais et s'arrête quand le nombre est trouvé. Ajouter une limite de 7 essais.}

\medskip
\textbf{Plan B si un élève termine tôt\,:} afficher un triangle d'étoiles avec deux boucles imbriquées.
\end{etape}

%------------------------------------------------------------------
\miseencommun{un élève présente sa boucle \texttt{for} et sa table de trace. Valoriser la prédiction avant l'exécution.}{un binôme montre son mot de passe à trois essais et explique la condition d'arrêt. Valoriser la précision.}{un élève fait jouer la classe à «~devine le nombre~». Valoriser la structure et l'interactivité.}{\textit{une boucle automatise une répétition : \texttt{for} pour un nombre de tours connu, \texttt{while} tant qu'une condition est vraie.}}

\cloture{Quelles valeurs prend \texttt{i} dans \texttt{for i in range(4):} ?}{Pourquoi la boucle \texttt{while n > 0:} ne s'arrête-t-elle jamais si on oublie \texttt{n = n - 1} ?}

%==================================================================
% ANNEXES
%==================================================================
\newpage
\section*{Annexe A : Support résumé de la mini-leçon}
\entete{blue}{Boucles en Python}

\begin{tcolorbox}[colback=yellow!10,colframe=yellow!50!black]
\textbf{for} : nombre de tours connu. \textbf{while} : tant que la condition est vraie. Dans les deux cas, le bloc répété est \textbf{indenté}.
\end{tcolorbox}

\begin{python}[for]
for i in range(1, 6):
    print(i)
\end{python}

\begin{python}[while]
essais = 0
while essais < 3:
    print("Essai", essais + 1)
    essais = essais + 1
\end{python}

\begin{tabularx}{\linewidth}{>{\ttfamily}p{3.2cm}Y}
\toprule
Instruction & Valeurs\\
\midrule
range(5) & 0, 1, 2, 3, 4\\
range(1, 6) & 1, 2, 3, 4, 5\\
range(0, 10, 2) & 0, 2, 4, 6, 8\\
\bottomrule
\end{tabularx}

%==================================================================
\newpage
\section*{Annexe B : Fiche de pratique guidée}
\entete{green}{Pratique guidée : répéter}

\subsection*{3a. Table de multiplication}
\begin{python}[À compléter pendant le codage en direct]
n = int(input("Quelle table ? "))
for i in range(1, _____):
    print(n, "x", i, "=", n * _____)
\end{python}

\subsection*{3b. Mot de passe}
\begin{python}[À compléter]
mot = input("Mot de passe : ")
while mot != "robot":
    mot = _____("Mot de passe : ")
print("Acces autorise")
\end{python}

\begin{tabularx}{\linewidth}{>{\bfseries}p{3.2cm}Y Y}
\toprule
Tour & Valeur de \texttt{mot} & La condition est vraie ?\\
\midrule
1 & & \\ 2 & & \\ 3 & & \\
\bottomrule
\end{tabularx}

\begin{tcolorbox}[colback=yellow!10,colframe=yellow!50!black,title={Cartes d'aide à utiliser dans l'ordre},fonttitle=\bfseries\small]
\textbf{Aide 1 :} regarde le programme du 3a et la table de trace de l'annexe A.\\
\textbf{Aide 2 :} Combien de fois la boucle se répète-t-elle ? Quelle variable change à chaque tour ?\\
\textbf{Aide 3 :} demande à l'enseignant d'écrire la ligne \texttt{while} avec toi.
\end{tcolorbox}

%==================================================================
\newpage
\section*{Annexe C : Niveau A (en dessous du niveau)}
\entete{orange}{Niveau A : je complète la boucle}

\textbf{Banque de mots :} for, in, range, print, 11, i.

\begin{python}[niveauA.py]
_____ i _____ range(1, _____):
    _____(i)
\end{python}

\subsection*{Activité : table de trace}
\begin{tabularx}{\linewidth}{>{\bfseries}p{3cm}Y}
\toprule
Tour & Valeur de \texttt{i}\\
\midrule
1 & \\ 2 & \\ 3 & \\ 10 & \\
\bottomrule
\end{tabularx}

\subsection*{Activité : vrai ou faux ?}
1. \texttt{range(5)} commence à 1. \cb Vrai \cb Faux\\
2. \texttt{range(1, 4)} donne 1, 2, 3. \cb Vrai \cb Faux

%==================================================================
\newpage
\section*{Annexe D : Niveau B (au niveau)}
\entete{blue}{Niveau B : trois programmes}

\subsection*{Programme 1 : table choisie}
Demande un nombre et affiche sa table de 1 à 10.

\subsection*{Programme 2 : somme de 1 à $n$}
Demande $n$ et affiche la somme $1 + 2 + \dots + n$. Indice : une variable \texttt{somme} qui part de 0.

\subsection*{Programme 3 : mot de passe à trois essais}
Le programme s'arrête après un bon mot de passe ou après trois essais ratés.

\begin{tabularx}{\linewidth}{>{\bfseries}p{3.2cm}Y Y Y}
\toprule
Programme & Entrée testée & Résultat attendu & Obtenu\\
\midrule
2 & $n = 4$ & & \\
2 & $n = 10$ & & \\
3 & bon mot dès le 2$^{\text{e}}$ essai & & \\
3 & trois erreurs & & \\
\bottomrule
\end{tabularx}

%==================================================================
\newpage
\section*{Annexe E : Niveau C (au-dessus du niveau)}
\entete{violet}{Niveau C : devine le nombre}

\begin{python}[niveauC.py : début du programme]
import random
secret = random.randint(1, 100)
essais = 0
trouve = False
while not trouve and essais < 7:
    # a toi de jouer : demander, comparer, compter
    pass
\end{python}

\begin{enumerate}
 \item Complète la boucle : demande un nombre, indique «~plus grand~» ou «~plus petit~», compte les essais.
 \item Affiche un message de fin différent en cas de victoire ou de défaite.
 \item Combien d'essais suffisent toujours avec la stratégie de la dichotomie (couper l'intervalle en deux) ? \dotfill
\end{enumerate}

\textbf{Défi bonus :} propose de rejouer à la fin avec une deuxième boucle.

%==================================================================
\newpage
\section*{Annexe F : Corrections et attendus (enseignant)}

\subsection*{Annexe B}
3a : \texttt{range(1, 11)}, \texttt{n * i}. 3b : \texttt{input}. La condition est vraie tant que le mot est différent de \texttt{"robot"}.

\subsection*{Annexe C}
\texttt{for i in range(1, 11): print(i)}. Trace : 1, 2, 3, 10. Vrai ou faux : Faux (commence à 0) ; Vrai.

\subsection*{Annexe D}
Somme : $n = 4$ donne 10 ; $n = 10$ donne 55. Mot de passe : \texttt{essais = 0} ; \texttt{while essais < 3 and mot != "robot"}. Cas limites : bon mot au 2\textsuperscript{e} essai, ou trois erreurs.

\subsection*{Annexe E}
Boucle attendue : \texttt{reponse = int(input(...))} ; \texttt{essais = essais + 1} ; \texttt{if reponse == secret: trouve = True} ; \texttt{elif reponse < secret: print("plus grand")} ; \texttt{else: print("plus petit")}. Dichotomie : 7 essais suffisent pour 1 à 100 ($2^7 = 128$).\\
\textbf{Critères de réussite :} boucle qui s'arrête, compteur correct, messages de fin distincts.

%==================================================================
\newpage
\section*{Annexe G : Exit ticket et grille d'observation}
\entete{gray}{Exit ticket : boucles}

\begin{enumerate}
 \item Quelles valeurs prend \texttt{i} dans \texttt{for i in range(4):} ?\\ \lig
 \item Pourquoi la boucle \texttt{while n > 0:} ne s'arrête-t-elle jamais si on oublie \texttt{n = n - 1} ?\\ \lig
 \item \autoeval
\end{enumerate}
\textbf{Attendus :} 1. 0, 1, 2, 3. 2. La condition reste toujours vraie, car \texttt{n} ne change pas.

\bigskip
\grille{\texttt{range} compris}{Condition d'arrêt}{Programme testé}

\end{document}
